\section{Materials and methods}
\label{sec:methods}

\subsection{Data, windows and the empirical target}
\label{sec:meth-data}

We used the National Healthcare Safety Network (NHSN) weekly hospital respiratory-data snapshot retrieved on 20 October 2026, restricting influenza admissions to the 50 US states and the District of Columbia \citep{cdc2026nhsn,frutos2025evaluation}. National and regional aggregates and territories were excluded. Hospital participation varied during the series, including the voluntary-reporting interval in May--October 2024 and the subsequent reporting requirement \citep{cms2024reporting}. The analysis is descriptive of this surveillance system and does not estimate an effect of that policy change.

The primary windows were the previously disclosed early-growth phases of four seasons: 1 October--26 November 2022, 7 October--23 December 2023, 9 November--28 December 2024 and 4 October--27 December 2025. Each comprises the contiguous national-growth segment with weekly growth at least 1.05 preceding the first turning point after the season's admissible start. The 2024 start follows resumed mandatory reporting. This rule and the primary statistic had been selected in earlier analysis rounds; the present revision is not a preregistered study. We preserved the windows rather than selecting them again to strengthen a fitted intercept. Inference conditions on their selection.

We rebuilt observations from the raw snapshot, preserving missing or suppressed counts as missing. A row requires four observed counts on consecutive weekly dates, from two weeks before its denominator week through the next week. Missing weeks are not treated as zero admissions. We retained genuine zero counts. The weighted scale must satisfy $D_t\ge5$; the primary single-week statistic also requires $C_t\ge5$. The resulting samples contain \Nra{} and \Ngs{} state-weeks for $R_A$ and $G$, respectively, across \Nstates{} jurisdictions and \Nweeks{} calendar weeks. The difference between samples is determined by these denominator conditions. State-week keys and original counts are frozen before any thinning analysis.

\subsection{Apparent statistics and relative squared deviation}
\label{sec:meth-stats}

Let $C_t$ denote admissions in week $t$. We used the weighted growth statistic
\begin{equation}
G_t=\frac{C_{t+1}}{D_t},\qquad D_t=0.65C_t+0.25C_{t-1}+0.10C_{t-2},
\label{eq:dt}
\end{equation}
and the growth-to-reproduction transform
\begin{equation}
R_{A,t}=\left[1+\frac{\theta}{7}\log\frac{C_{t+1}+0.5}{C_t+0.5}\right]_{+}^{\alpha},\qquad
\alpha=(g_\mu/g_\sigma)^2,\quad\theta=g_\sigma^2/g_\mu.
\label{eq:ra}
\end{equation}
The positive-part convention defines the transformation when its base is nonpositive. It is a numerical convention, not an assertion that the exponential-growth approximation holds there. The gamma-interval mapping follows the renewal relationship between growth and reproduction \citep{wallinga2007how}. Neither statistic is an observed true reproduction number or a posterior estimate from an infection-renewal model \citep{cori2013new,eales2024differences}.

For continuity, the primary transform uses $g_\mu=3.6$ days and $g_\sigma=1.6$ days, the household serial-interval estimate of \citet{cowling2009estimation}. We distinguish this proxy from a generation interval. Sensitivities use mean 3.2 days with the original coefficient of variation held fixed, and mean 3.2 days with SD 2.1 days, the intrinsic generation-time estimates in Table 2 of \citet{chan2025estimating}. The fixed-coefficient-of-variation choice is an assumption. The reported household realised generation time and subtype serial intervals are different quantities \citep{chan2025estimating,vink2014serial}.

For each specification, we calculated
\begin{equation}
L_{st}=\left(\frac{X_{st}-\bar X}{\bar X}\right)^2,
\qquad E(L_{st}\mid D_{st})=a+b/D_{st},
\label{eq:affine}
\end{equation}
where $X$ is an apparent statistic and $\bar X$ is its pooled mean. For $R_A$, $C_t$ replaces $D_t$ in the regression. This response measures dispersion around a fitted reference; it is not squared estimation error against known $R_t$. The ratio $b/a$, when defined, describes the affine working model on that statistic's count scale. It is not directly the latent balance scale $\mx$.

\subsection{Fitting, uncertainty and model comparison}
\label{sec:meth-inference}

The affine mean was fitted by iterated feasible generalised least squares (FGLS), with precision proportional to the inverse squared fitted mean. Intercepts and slopes can be signed, while fitted means on observed rows must remain positive. Damped updates decrease the gamma quasi-deviance objective $\sum_iw_i[L_i/\mu_i+\log\mu_i]$; a positive-mean numerical fallback uses the same objective. We record convergence, rank and failures. Reaching an iteration limit does not produce a reported estimate. This objective is an estimating criterion, not a claim that the empirical losses follow a gamma likelihood.

Primary pooled intervals use 2000 state $\times$ calendar-week pigeonhole bootstrap draws \citep{owen2007pigeonhole}. A row receives the product of its independently drawn state and week multiplicities. The reference mean and regression are re-estimated in each draw. Secondary analyses use 1000 draws. We also resample two- and three-week blocks within seasons, combined with state multiplicities, to examine sensitivity to persistence in national shocks. These few seasons do not support unrestricted inference across arbitrary future epidemic regimes. Undefined ratios, denominator signs and failed fits are retained in diagnostics; any finite-draw interval is labelled accordingly.

We compared the affine mean with $\lambda D^{-\gamma}$ by leaving out one state or one season at a time. All response centres are estimated from the training observations; held-out losses use that training centre. Each model is fitted only on its training fold. Held-out scores are the mean of $L/\mu+\log\mu$, with squared prediction errors retained as a secondary measure. We report power-law minus affine scores, so a positive difference favours the affine model. The overlapping training folds are not independent replications, and four season folds are not used to manufacture a small-sample significance test. The earlier residual-bootstrap platform $p$-value is not carried forward as a corrected test.

\subsection{Controlled binomial thinning}
\label{sec:meth-thinning}

At $q\in\{0.8,0.6,0.5,0.4,0.3,0.2,0.1\}$, each unique state-week count is sampled once from $\operatorname{Bin}(C,q)$ per realisation and reused in every overlapping ratio. No eligibility filter or window is reapplied. We recompute the statistic, its centre and the affine fit using the original $C_t$ as regressor scale. The point response averages \ThinPointReps{} separately fitted realisations. Each of \ThinBootstrap{} outer bootstrap draws averages \ThinInnerReps{} inner fits per level. Inner counts are increased uniformly when the estimated Monte Carlo standard error exceeds the prespecified fraction of sample uncertainty; actual budgets and remaining Monte Carlo error are reported.

The approximate relation is
\begin{equation}
b(q)-b(1)\simeq c_{\rm thin}(1/q-1).
\label{eq:thinning}
\end{equation}
We fit the through-origin rate, retain each level and repeat the summary without $q=0.1$ and with each level omitted in turn. Prespecified subsets require original $C_t\ge50$, or both $C_t,C_{t+1}\ge50$. These conditions bound expected thinned counts, not every random realisation. A reference coefficient is the weighted mean of $\kappa(g)^2[1+(C_t+0.5)/(C_{t+1}+0.5)]$, where $\kappa(g)=(g_\mu/7)/[1+\theta\log(g)/7]$ and $g=(C_{t+1}+0.5)/(C_t+0.5)$. It approximates the relative transformation response and stabilises inverse counts; it is not the exact variance of the recentered nonlinear statistic. We recompute this coefficient in every bootstrap draw to obtain paired $c_{\rm thin}/b_P$ intervals. Simultaneous bands describe the vector of intercept changes. A known Poisson-count example runs through the same procedure as a mechanism check, not as a coverage validation.

\subsection{State heterogeneity and reference sensitivity}
\label{sec:meth-fe}

With the same pooled-loss definition, we fit $E(L_{st}\mid C_{st})=a_s+b_{\rm w}/C_{st}$. State summaries target sampled state occurrences: repeated draws of a state count repeatedly, and absent states do not enter as zero coefficients. Row- and case-weighted summaries use resampled totals. Pooled and within-state contrasts are paired within draws. We separately replace $C_t$ by the actual weighted $D_t$. Week-only and state-plus-week fits are numerical diagnostics; failed fits are explicitly unavailable.

A separate analysis changes the reference in the loss, using season means, season-specific linear trends common to states, national-week means, and leave-own-state-out national-week means. The global mean squared remains the denominator, isolating centring from renormalisation. All centres and trends are refitted under bootstrap weights. These operations change the empirical target and may remove real transmission signal. Their fitted intercept differences are therefore specification sensitivities, not orthogonal biological variance components. We additionally report $\operatorname{Var}_w(\bar d_w)/\operatorname{Var}_{st}(d_{st})$ for signed deviations $d=X/\bar X-1$, explicitly as a descriptive variance ratio rather than an ICC of squared losses.

\subsection{Conditional calibration and design illustrations}
\label{sec:meth-inversion}

A daily renewal simulator reproduces the four fixed window lengths and initial state counts. Transmission has independent negative-binomial offspring; weekly reproduction-number shocks have separate national and state components with declared marginal variance and persistence. Admission delays and report probabilities are assumptions unless documented externally. Binomial admissions are allocated over delays. An alternative observation model samples directly from the latent delayed admission intensity $\lambda$, with conditional mean $\lambda$ and variance $\varphi\lambda$; it does not resample an already noisy admission count. The Poisson limit is $\varphi=1$. Observation-moment tests check these definitions independently.

We apply the same statistics, filters and fitting functions to simulated panels. The variance parameter describes weekly reproduction shocks within this simulator. It is not converted to a per-generation variance or an operational weekly case threshold. Paired empirical-intercept draws and paired simulated-panel resampling propagate sampling and calibration Monte Carlo uncertainty. Separate generation-interval and observation-model scenarios describe structural sensitivity; they are not a posterior distribution or a joint confidence interval over mechanisms. Supplementary methods specify all settings, boundaries and simulation diagnostics. Further supplementary analyses report season and participation sensitivities and an explicitly hypothetical curve-discrimination design experiment. None is used to infer a causal reporting-mandate effect or a universal surveillance sample-size requirement.
